\documentclass[11pt,a4paper]{article}
\usepackage{amsmath,amssymb,amsfonts,amsthm}
\usepackage{geometry}
\geometry{margin=1in}
\title{\textbf{Geotemporal Hydrodynamics (GTH) v12.0\\Paper 3: Force-Free Beltrami Vortex Wakes}}
\author{CoderQuan2}
\date{\today}
\begin{document}
\maketitle
\begin{abstract}
We derive the structure of self-stabilizing, force-free Beltrami vortex wakes in continuous GTH flow fields, demonstrating that hydrodynamic helicity conservation ensures persistent structural stability.
\end{abstract}
\section{Beltrami Field Dynamics}
A velocity field $\mathbf{u}$ is defined as a Beltrami field if it aligns with its own curl:
\begin{equation}
\nabla \times \mathbf{u} = \lambda \mathbf{u}
\end{equation}
The force-free condition eliminates advective non-linear momentum transport loss:
\begin{equation}
(\mathbf{u} \cdot \nabla) \mathbf{u} = \nabla \left(\frac{|\mathbf{u}|^2}{2}\right) - \mathbf{u} \times (\nabla \times \mathbf{u}) = \nabla \left(\frac{|\mathbf{u}|^2}{2}\right)
\end{equation}
\end{document}
